\chapter[Δεδομένα \& Χαρακτηριστικά]{Δεδομένα και Μηχανική Χαρακτηριστικών}
\label{ch:data}

Η ποιότητα, η πληρότητα και η επεξεργασία των δεδομένων καθορίζουν την
απόδοση κάθε μοντέλου μηχανικής μάθησης, ιδιαίτερα στην EPF όπου οι χρονοσειρές εμφανίζουν
έντονη μη γραμμικότητα, πολλαπλή εποχικότητα και ακραίες τιμές (price spikes)~\cite{kaufman2012leakage}.
Στην παρούσα εργασία αντλούνται δεδομένα από οκτώ ξεχωριστές πηγές, που καλύπτουν την
περίοδο Ιανουάριος 2017 - Μάρτιος 2026 με ωριαία ανάλυση. Η ανομοιογένεια των πηγών
(διαφορετικές χωρικές κλίμακες, ημερήσια αντί ωριαία ανάλυση, διαφορετικές ζώνες ώρας,
κενές τιμές) απαιτεί ενιαία διοχέτευση κανονικοποίησης και συγχώνευσης πριν την κατασκευή
χαρακτηριστικών.

Το παρόν κεφάλαιο αναλύει διαδοχικά: (α) τις πηγές ακατέργαστων δεδομένων και τη δομή
τους, (β) τον αγωγό κανονικοποίησης και συγχώνευσης, (γ) τη μηχανική χαρακτηριστικών,
(δ) την αντιμετώπιση διαρροής δεδομένων, (ε) τον τελικό πίνακα χαρακτηριστικών και
(στ) τον διαχωρισμό εκπαίδευσης/αξιολόγησης.


\section{Πηγές Ακατέργαστων Δεδομένων}
\label{sec:data_sources}

Ο Πίνακας~\ref{tab:data_sources} συνοψίζει τις οκτώ πηγές δεδομένων που
χρησιμοποιούνται στην εργασία, με αναφορά στην κάθε πηγή και τον σύνδεσμο
πρόσβασης.

\begin{table}[htbp]
\centering
\caption{Πηγές ακατέργαστων δεδομένων.}
\label{tab:data_sources}
\small
\begin{tabular}{p{3.8cm}p{2cm}p{3.2cm}p{4.1cm}}
\toprule
Δεδομένα & Ανάλυση & Πηγή & Σύνδεσμος \\
\midrule
Τιμές DAM (€/MWh) & Ωριαία & ENTSO-E Transparency Platform~\cite{entsoe_transparency} & \url{transparency.entsoe.eu} \\
Πραγματικό Φορτίο (MW) & Ωριαία & ENTSO-E Transparency Platform~\cite{entsoe_transparency} & \url{transparency.entsoe.eu} \\
Παραγωγή ανά Τεχνολογία (MW) & 15\,λεπτά & ENTSO-E Transparency Platform~\cite{entsoe_transparency} & \url{transparency.entsoe.eu} \\
Τιμή Φυσικού Αερίου TTF (€/MWh) & Ημερήσια & Investing.com~\cite{ttf_gas} & \url{investing.com} \\
Τιμή CO\textsubscript{2} EUA (€/τ) & Ημερήσια & Investing.com~\cite{co2_ets} & \url{investing.com} \\
Προβλέψεις ΑΠΕ day-ahead & Ωριαία & ENTSO-E Transparency Platform~\cite{entsoe_transparency} & \url{transparency.entsoe.eu} \\
Μετεωρολογικά (6 τοποθεσίες) & Ωριαία & Open-Meteo API~\cite{openmeteo} & \url{open-meteo.com/api} \\
Πρόβλεψη Φορτίου ΑΔΜΗΕ & Ωριαία & ΑΔΜΗΕ / proxy \texttt{load\_lag24} & \url{admie.gr} \\
\bottomrule
\end{tabular}
\end{table}

\subsection{Τιμές Αγοράς Επόμενης Ημέρας — ENTSO-E DAM}

Οι ωριαίες τιμές εκκαθάρισης της αγοράς επόμενης ημέρας (DAM) για τη Ζώνη Τιμολόγησης
Ελλάδας (BZN|GR) αντλούνται από την πλατφόρμα ENTSO-E Transparency
Platform~\cite{entsoe_transparency}. Τα ακατέργαστα δεδομένα αποθηκεύονται τοπικά σε 11
αρχεία CSV της μορφής \texttt{GUI\_ENERGY\_PRICES\_YYYYMMDD-YYYYMMDD.csv}, που
καλύπτουν το διάστημα Δεκέμβριος 2014 - Μάρτιος 2026. Κάθε εγγραφή αντιστοιχεί σε μία
ώρα λειτουργίας και εκφράζεται σε €/MWh.

Από το 2024 και εξής, η ENTSO-E μετάβηκε σε ανάλυση 15 λεπτών για την Ελλάδα, αποδίδοντας
4 τιμές ανά ώρα αντί μίας. Στο πλαίσιο της παρούσας εργασίας, η χρονική ανάλυση ορίστηκε
ωριαία σε ολόκληρο τον ορίζοντα μελέτης. Η μετατροπή γίνεται αυτόματα: οι χρονοσφραγίδες
στρογγυλοποιούνται στην πλησιέστερη ώρα (\texttt{floor("H")}) και οι 4 τετάρτωρες τιμές
συγχωνεύονται σε μία ωριαία τιμή μέσω αριθμητικού μέσου. Για τα παλαιότερα δεδομένα
ωριαίας ανάλυσης, η λειτουργία αυτή παραμένει διαφανής. Η ίδια λογική εφαρμόζεται και στα
δεδομένα φορτίου. Η σειρά τιμών είναι ο στόχος πρόβλεψης $y_t$ για την εργασία EPF.


\subsection{Πραγματικό Ηλεκτρικό Φορτίο — ENTSO-E}

Τα ωριαία δεδομένα πραγματικής κατανάλωσης ηλεκτρικής ενέργειας (Actual Total Load) για
την Ελλάδα σε MW αντλούνται επίσης από το ENTSO-E Transparency
Platform~\cite{entsoe_transparency}. Τα ακατέργαστα αρχεία CSV βρίσκονται στον κατάλογο
\texttt{data/raw/load/} και καλύπτουν Ιανουάριο 2015 - Μάρτιο 2026. Εφαρμόζεται η ίδια
λογική χειρισμού δεδομένων 15 λεπτών (\texttt{mean()}) με αυτήν της σειράς τιμών.

Η σειρά φορτίου έχει διπλό ρόλο: αφενός είναι ο στόχος πρόβλεψης $y_t$ για
την εργασία STLF, αφετέρου χρησιμεύει ως εξωγενές χαρακτηριστικό στα μοντέλα τιμής, υπό
τη μορφή χρονικά καθυστερημένων τιμών (βλ.\ Ενότητα~\ref{sec:lag_features}).


\subsection{Πραγματική Παραγωγή ανά Τύπο Τεχνολογίας — ENTSO-E}

Τα δεδομένα παραγωγής ηλεκτρικής ενέργειας ανά τύπο τεχνολογίας σε MW αντλούνται από
το ENTSO-E Transparency Platform~\cite{entsoe_transparency}. Τα αρχεία
(\texttt{AGGREGATED\_GENERATION\_PER\_TYPE\_*.csv}) καλύπτουν Δεκέμβριο 2014 -
Μάρτιο 2026 σε ανάλυση 15 λεπτών. Εντός της διαδικασίας επεξεργασίας, κάθε αρχείο CSV
αναπτύσσεται σε πίνακα με τη χρήση \texttt{pivot\_table} (index: χρονική σφραγίδα,
columns: τύπος τεχνολογίας, aggfunc: ``sum''), παράγοντας πέντε στήλες:
\texttt{gen\_solar}, \texttt{gen\_wind}, \texttt{gen\_gas}, \texttt{gen\_lignite},
\texttt{gen\_hydro}. Τα δεδομένα αυτά αντιπροσωπεύουν πραγματοποιηθείσα παραγωγή που
δεν είναι γνωστή τη στιγμή υποβολής προσφορών (D-1, \textasciitilde12:00 CET) και
χρησιμοποιούνται αποκλειστικά υπό χρονική καθυστέρηση (βλ.\ Ενότητα~\ref{sec:leakage}).


\subsection{Τιμή Φυσικού Αερίου — Dutch TTF Futures (Ημερήσια)}

Η ημερήσια τιμή φυσικού αερίου αντλείται από τα συμβόλαια μελλοντικής εκπλήρωσης του
δείκτη Dutch TTF (Title Transfer Facility), που αποτελεί τον κυρίαρχο ευρωπαϊκό δείκτη
αναφοράς για τη χονδρική αγορά αερίου~\cite{ttf_gas}. Τα δεδομένα λήφθηκαν χειροκίνητα
από την πλατφόρμα Investing.com, εκφράζονται σε €/MWh και χρησιμοποιούνται αποκλειστικά
στην εργασία τιμής. Κατά τη συγχώνευση, οι ημερήσιες τιμές επεκτείνονται με \texttt{ffill()}
σε ωριαία ανάλυση.

Δεδομένου ότι η παραγωγή από φυσικό αέριο συχνά καθορίζει οριακά την τιμή εκκαθάρισης
στην ελληνική αγορά, η τιμή TTF ασκεί ισχυρή αιτιακή επίδραση στη DAM~\cite{lago2021}.
Για το διάστημα Νοέμβριος 2025 - Φεβρουάριος 2026 χρησιμοποιήθηκαν συνθετικές τιμές
με γραμμική παρεμβολή μεταξύ σημείων αναφοράς (Δεκ.\ 2025: 46 €/MWh, Ιαν.\ 21, 2026:
36,77 €/MWh, Φεβ.\ 20, 2026: 35 €/MWh). Επειδή η παρεμβολή βασίζεται σε
μεταγενέστερα σημεία αναφοράς, εισάγει ήπια εξάρτηση από το μέλλον (look-ahead)
στο χαρακτηριστικό gas\_price για το διάστημα δοκιμής· ο περιορισμός αυτός και η
ενδεδειγμένη αντιμετώπισή του (αντικατάσταση με πραγματικά δεδομένα TTF)
συζητούνται στο Κεφάλαιο~\ref{ch:conclusions}.


\subsection{Τιμή Δικαιωμάτων CO\textsubscript{2} — EU ETS Futures (Ημερήσια)}

Οι ημερήσιες τιμές δικαιωμάτων εκπομπών CO$_2$ αντλήθηκαν από την πλατφόρμα
Investing.com και αντιστοιχούν στα συμβόλαια EUA (European Union Allowance) του EU
ETS, εκφρασμένες σε €/τόνο CO$_2$~\cite{co2_ets}. Το κόστος CO$_2$ ενσωματώνεται ως
μεταβλητό κόστος παραγωγής για κάθε μονάδα που χρησιμοποιεί ορυκτά καύσιμα και
χρησιμοποιείται αποκλειστικά για \texttt{task=price}. Η επέκτασή του σε ωριαία ανάλυση
ακολουθεί το ίδιο σχήμα \texttt{ffill()} με την τιμή αερίου. Η χρήση σύγχρονων ημερήσιων
τιμών καυσίμων αποτελεί οριακή επιλογή από πλευρά leakage, αποδεκτή στη βιβλιογραφία
EPF~\cite{lago2021}.


\subsection{Προβλέψεις Παραγωγής Επόμενης Ημέρας — ENTSO-E Extra}

Ωριαίες προβλέψεις παραγωγής από ΑΠΕ για την επόμενη ημέρα (day-ahead) (\path{solar/wind/gen_fc_dayahead}), έτσι
όπως δημοσιεύονται από τη ENTSO-E πριν τη διεξαγωγή της DAM, φορτώνονται από
\path{data/processed/entsoe_extra_hourly.parquet} (71.640 παρατηρήσεις,
Οκτ.\ 2017 - Δεκ.\ 2025). Τα δεδομένα αυτά αποτελούν νόμιμα σύγχρονα χαρακτηριστικά:
δημοσιεύονται πριν την υποβολή προσφορών. Για κάθε στήλη δημιουργείται δυαδικός δείκτης
απόντος (\texttt{*\_missing}), παράγοντας 6 συνολικές στήλες.


\subsection{Μετεωρολογικά Δεδομένα — Open-Meteo API}

Ωριαία μετεωρολογικά δεδομένα αντλούνται από τη δωρεάν ιστορική API της υπηρεσίας
Open-Meteo~\cite{openmeteo}. Τα επεξεργασμένα δεδομένα αποθηκεύονται στο
\texttt{data/processed/weather\_gr\_hourly.parquet} (73.416 παρατηρήσεις,
Οκτ.\ 2017 - Μαρ.\ 2026). Για κάθε μία από τις έξι γεωγραφικές τοποθεσίες
(\texttt{gr\_mean}, Αθήνα, Θεσσαλονίκη, Πάτρα, Λάρισα, Ηράκλειο) αντλούνται επτά
μετεωρολογικές μεταβλητές (θερμοκρασία, υγρασία, κατακρημνίσματα, νεφοκάλυψη, ταχύτητα
ανέμου, ριπές ανέμου, ηλιακή ακτινοβολία), αποδίδοντας $6 \times 7 = 42$ πρωτογενείς στήλες.
Για κάθε στήλη δημιουργείται δυαδικός δείκτης απόντος, οδηγώντας σε 84 συνολικές στήλες.
Τα κενά συμπληρώνονται σε δύο βήματα: \texttt{ffill(limit=6)} και ακολούθως συμπλήρωση με μηδενικά (zero-fill). Η συμπλήρωση με μηδενικά είναι συνήθης πρακτική όταν συνοδεύεται από δυαδικούς δείκτες απουσίας: τα δενδρικά μοντέλα μπορούν τότε να διακρίνουν μια πραγματική μηδενική τιμή από μια ελλείπουσα, χωρίς να εισάγεται μεροληψία από τεκμαρτές τιμές.


\subsection{Πρόβλεψη Φορτίου Επόμενης Ημέρας (load\_fc)}

Η ωριαία day-ahead πρόβλεψη φορτίου που δημοσιεύει ο ΑΔΜΗΕ αποτελεί νόμιμο σύγχρονο
χαρακτηριστικό αποκλειστικά για \texttt{task=price}. Το αρχείο
{data/processed/load\_forecast\_hourly.parquet} (Οκτ.\ 2025 - Δεκ.\ 2025, 2.207 ώρες)
δεν καλύπτει ολόκληρο το εκπαιδευτικό εύρος· για αυτό εφαρμόζεται υποκατάστατο (proxy) με \texttt{load\_lag24} (πραγματικό φορτίο ίδιας ώρας 24 ώρες νωρίτερα). Αξιολόγηση έδειξε
ότι η αντικατάσταση ακριβέστερης πρόβλεψης (MAE = 80,7 MW) με την ENTSO-E πρόβλεψη
(MAE = 267 MW) μεταβάλλει το τελικό MAE τιμής κατά $<0{,}1\%$, υποδηλώνοντας ότι τα
μοντέλα τιμής βασίζονται κυρίως στις χρονικές υστερήσεις τιμής και στα επίπεδα καυσίμων.


\section{Κανονικοποίηση Χρονικού Δείκτη}
\label{sec:time_normalization}

Μία από τις σημαντικότερες προκλήσεις κατά τη συγχώνευση ανομοιογενών πηγών είναι η
αντιμετώπιση ανακολουθιών στη χρονική ανάλυση και τη ζώνη ώρας. Η διοχέτευση εφαρμόζει δύο
βοηθητικές συναρτήσεις πριν από κάθε join.

Η \texttt{\_ensure\_hourly\_index()} μετατρέπει τον δείκτη κάθε σειράς δεδομένων σε
\texttt{pd.DatetimeIndex}, εφαρμόζει \texttt{floor("H")}, ταξινομεί χρονολογικά και αφαιρεί
τιμές NaT. Εάν υπάρχουν πολλαπλές εγγραφές για την ίδια ώρα (π.χ.\ δεδομένα 15 λεπτών),
εφαρμόζεται \texttt{groupby(level=0).mean()}, ο οποίος υπολογίζει τον αριθμητικό μέσο όρο
των τιμών, εξασφαλίζοντας ότι η τελική εγγραφή κάθε ώρας αντιπροσωπεύει αξιόπιστα τις
πραγματικές συνθήκες.

Η \texttt{\_to\_utc\_naive\_index()} αντιμετωπίζει τις ζώνες ώρας: μετατρέπεται πρώτα σε
UTC και στη συνέχεια αφαιρείται η πληροφορία ζώνης (UTC-naive), ώστε να αποτραπούν
διπλές ή κενές εγγραφές κατά τις αλλαγές θερινής/χειμερινής ώρας. Κατά τη συγχώνευση,
ο αγωγός ανιχνεύει αυτόματα πιθανή μετατόπιση ζώνης ώρας, δοκιμάζοντας
$\Delta \in \{-3,\ldots,+3\}$ ωρών.


\section{Διοχέτευση Επεξεργασίας Δεδομένων}
\label{sec:pipeline}

Η μετατροπή των ακατέργαστων πηγών στον τελικό πίνακα χαρακτηριστικών γίνεται σε πέντε
διαδοχικές φάσεις, που συνοψίζονται γραφικά στο Σχήμα~\ref{fig:pipeline} (Ενότητα~\ref{subsec:final_table}).

Η κεντρική συνάρτηση \texttt{process\_hourly()} του \texttt{src/data.py} διαχειρίζεται
ολόκληρη τη διοχέτευση, εκτελούμενη ξεχωριστά για \texttt{task=price} και \texttt{task=load}.
Οι δύο εκδοχές μοιράζονται τον ίδιο πυρήνα επεξεργασίας αλλά διαφέρουν στις εξωγενείς
μεταβλητές: η εκδοχή \texttt{task=price} ενσωματώνει τις μεταβλητές καυσίμων και την
πρόβλεψη φορτίου, τις οποίες η \texttt{task=load} παραλείπει εξ ορισμού.


\section{Μηχανική Χαρακτηριστικών}
\label{sec:feature_engineering}

Ο πίνακας χαρακτηριστικών κατασκευάζεται αποκλειστικά από πληροφορίες γνωστές κατά τον
πραγματικό χρόνο πρόβλεψης (D-1). Κάθε παρατήρηση $(x_t, y_t)$ αντιστοιχεί σε μία ώρα
$t$, με $x_t \in \mathbb{R}^d$ και $y_t$ την τιμή-στόχο σε €/MWh ή MW.


\subsection{Ημερολογιακά Χαρακτηριστικά}

Τα ημερολογιακά χαρακτηριστικά αποτυπώνουν γνωστές εκ των προτέρων δομές εποχικότητας
και ανθρώπινης δραστηριότητας~\cite{uniejewski2019}. Η ώρα ημέρας $h \in \{0,\ldots,23\}$
και η ημέρα εβδομάδας $d \in \{0,\ldots,6\}$ κωδικοποιούνται κυκλικά:

\begin{equation}
  \left(\sin\tfrac{2\pi h}{24},\; \cos\tfrac{2\pi h}{24},\;
        \sin\tfrac{2\pi d}{7},\; \cos\tfrac{2\pi d}{7}\right),
  \label{eq:cyclic}
\end{equation}

εξασφαλίζοντας γεωμετρική συνέχεια μεταξύ ώρας 23 και ώρας 0. Επιπλέον
συμπεριλαμβάνεται δυαδική μεταβλητή ελληνικών αργιών (\texttt{is\_holiday}) από τη
βιβλιοθήκη \texttt{holidays.Greece()}. Η ομάδα αυτή αποδίδει 7 χαρακτηριστικά, κοινά
για τις δύο εργασίες.


\subsection{Χρονικά Καθυστερημένα Χαρακτηριστικά}
\label{sec:lag_features}

Ο κύριος κορμός του πίνακα χαρακτηριστικών αποτελείται από χρονικά καθυστερημένες τιμές.
Το κεντρικό σύνολο υστερήσεων που εφαρμόζεται σε όλες τις κύριες χρονοσειρές είναι:

\begin{equation}
  L = \{1, 2, 3, 6, 12, 24, 48, 168\} \quad \text{(ώρες)},
  \label{eq:lag_set}
\end{equation}

επιλεγμένο ώστε να αποτυπωθεί η βραχυπρόθεσμη δυναμική (lags 1-3), οι ιδιαιτερότητες
αιχμής/εκτός αιχμής (lags 6, 12), η ημερήσια περιοδικότητα (lags 24, 48) και η εβδομαδιαία
εποχικότητα (lag 168). Η τιμή $y_{t-168}$ είναι ισχυρό αναλλοίωτο χαρακτηριστικό στη
βιβλιογραφία EPF~\cite{uniejewski2019}.

\begin{table}[htbp]
\centering
\caption{Σύνολα υστερήσεων ανά σήμα.}
\label{tab:lag_sets}
\begin{tabular}{llll}
\toprule
Σήμα & Περιγραφή & Σύνολο & \# lags \\
\midrule
$y$            & Τιμή DAM ή φορτίο (στόχος)   & $L$                      & 8 \\
load           & Πραγματικό φορτίο             & $L$                      & 8 \\
residual\_load & Φορτίο μείον ΑΠΕ παραγωγή     & $L$                      & 8 \\
gas\_price     & Τιμή φυσικού αερίου TTF        & $L$                      & 8 \\
co2\_price     & Τιμή δικαιωμάτων CO$_2$       & $L$                      & 8 \\
gen\_solar     & Παραγωγή ηλιακή               & $L' = \{1,2,24,48,168\}$ & 5 \\
gen\_wind      & Παραγωγή αιολική              & $L' = \{1,2,24,48,168\}$ & 5 \\
\bottomrule
\end{tabular}
\end{table}

Η υπολειμματική ζήτηση ορίζεται ως
$\text{residual\_load}_t = \text{load}_t - \text{gen\_solar}_t - \text{gen\_wind}_t$
και αποτυπώνει το τμήμα κατανάλωσης που πρέπει να καλυφθεί από συμβατικές μονάδες,
ως ένδειξη οριακής ζήτησης.


\subsection{Κυλιόμενα Στατιστικά}

Για τον στόχο $y$ υπολογίζονται δύο κυλιόμενοι μέσοι:

\begin{equation}
  y_{\text{roll24},t} = \frac{1}{24}\sum_{k=1}^{24} y_{t-k},
  \label{eq:roll24}
\end{equation}

\begin{equation}
  y_{\text{roll168},t} = \frac{1}{168}\sum_{k=1}^{168} y_{t-k}.
  \label{eq:roll168}
\end{equation}

Ο χρόνος ολίσθησης ξεκινά από $k=1$ ώστε να μην περιλαμβάνεται η τρέχουσα τιμή $y_t$,
αποτρέποντας διαρροή. Το $y_{\text{roll24}}$ λειτουργεί ως μέτρο βραχυπρόθεσμης τάσης,
ενώ το $y_{\text{roll168}}$ καταγράφει το εβδομαδιαίο εποχικό επίπεδο.


\subsection{Δείκτες Απόντων Τιμών}

Για κάθε στήλη που ενδέχεται να εμφανίζει κενές τιμές (42 μετεωρολογικές + 3 ENTSO-E
day-ahead), δημιουργείται δυαδική στήλη \texttt{\{col\}\_missing} που λαμβάνει τιμή 1
εάν η αρχική τιμή απουσίαζε. Η τεχνική αυτή επιτρέπει στα μοντέλα να διακρίνουν μεταξύ
«η τιμή είναι μηδέν» και «η τιμή είναι άγνωστη»~\cite{goodfellow2016deep}.


\subsection{Αντιμετώπιση Διαρροής Δεδομένων}
\label{sec:leakage}

Η διαρροή δεδομένων (data leakage)~\cite{kaufman2012leakage} είναι ένα από τα πιο συχνά και σοβαρά μεθοδολογικά σφάλματα στην εφαρμοσμένη πρόβλεψη χρονοσειρών και οδηγεί σε
αισιόδοξες αλλά μη αναπαραγώγιμες εκτιμήσεις απόδοσης. Στην παρούσα εργασία
εφαρμόζεται ρητή διόρθωση leakage σε δύο επίπεδα.

\textbf{Εργασία τιμής.} Κατά τη στιγμή υποβολής προσφορών (D-1, \textasciitilde12:00 CET),
οι ποσότητες της επόμενης ημέρας δεν έχουν πραγματοποιηθεί. Ο αγωγός αφαιρεί τις
σύγχρονες ($t=0$) τιμές των: load, residual\_load, gen\_solar, gen\_wind, gen\_gas,
gen\_lignite, gen\_hydro. Οι χρονικά καθυστερημένες εκδοχές ($k \geq 1$) παραμένουν.

\textbf{Εργασία φορτίου.} Αφαιρούνται οι σύγχρονες τιμές των: price, residual\_load,
gen\_solar, gen\_wind, gen\_gas, gen\_lignite, gen\_hydro.

\textbf{Cross-Validation χωρίς διαρροή.} Κατά τη βελτιστοποίηση υπερπαραμέτρων με
Optuna, το cross-validation εκτελείται με rolling-origin σχήμα: ο επικυρωτής πάντα
ακολουθεί χρονικά~\cite{bergmeir2018}, αποτρέποντας τη «διαρροή μέλλοντος».


\subsection{Πυκνές Ενδοημερήσιες Υστερήσεις (Dense Intraday Lags) — Μόνο για MIMO/Direct}

Το βασικό σύνολο $L$ δεν περιλαμβάνει τις ενδοημερήσιες υστερήσεις lag-4 έως lag-23 (εκτός των ήδη υπαρχόντων lag-6 και lag-12), δηλαδή 18 επιπλέον υστερήσεις. Οι στρατηγικές MIMO και Direct προβλέπουν ολόκληρο το 24ωρο ταυτόχρονα, ως πολλαπλή έξοδο και από μία μόνο αφετηρία, χωρίς αναδρομική τροφοδότηση· έτσι όλες οι ενδοημερήσιες τιμές $y_{t-4}, \ldots, y_{t-23}$ είναι ήδη γνωστές, παρατηρημένες τιμές τη στιγμή της πρόβλεψης. Επειδή συνιστούν καθαρή ιστορική πληροφορία και όχι προβλέψεις, προστίθενται ως \emph{πυκνές υστερήσεις} (dense lags) και εμπλουτίζουν την περιγραφή της ενδοημερήσιας δυναμικής:

\begin{itemize}
  \item Βασικό σύνολο (TF/Rec): 152 χαρακτηριστικά (τιμή) / 133 (φορτίο).
  \item Dense σύνολο (MIMO/Direct): $152 + 18 = 170$ (τιμή) / $133 + 18 = 151$ (φορτίο).
\end{itemize}

Στα MIMO/Direct οι dense lags παρέχουν αυτή την καθαρή ενδοημερήσια πληροφορία χωρίς συσσωρευμένο σφάλμα. Στη Recursive στρατηγική, αντίθετα, οι ίδιες υστερήσεις θα συμπληρώνονταν από προβλέψεις του ίδιου του μοντέλου και θα μετέφεραν συσσωρευμένο σφάλμα· γι' αυτό η προσθήκη τους επιδεινώνει την απόδοση (βλ.\ Κεφάλαιο~\ref{ch:experiments}).


\subsection{Τελικός Πίνακας Χαρακτηριστικών}
\label{subsec:final_table}

Ο Πίνακας~\ref{tab:features_summary} συνοψίζει τις ομάδες χαρακτηριστικών ανά εργασία.
Μετά το \texttt{dropna()}, τα επεξεργασμένα αρχεία αποθηκεύονται:

\begin{table}[htbp]
\centering
\caption{Επεξεργασμένα αρχεία Parquet μετά τη φάση 5 του αγωγού.}
\label{tab:parquet_files}
\small
\begin{tabular}{lllll}
\toprule
Αρχείο & Εργασία & Γραμμές & Στήλες & Κάλυψη \\
\midrule
\texttt{hourly.parquet}       & τιμή   & 80.126 & 153 & 27/01/2017 - 20/03/2026 \\
\texttt{hourly\_load.parquet} & φορτίο & 98.054 & 134 & 08/01/2015 - 19/03/2026 \\
\bottomrule
\end{tabular}

\smallskip
{\footnotesize Η στήλη «Στήλες» μετρά το σύνολο των στηλών του αρχείου,
συμπεριλαμβανομένου του στόχου $y$. Η \texttt{make\_xy()} αφαιρεί τον στόχο,
οπότε τα διαθέσιμα χαρακτηριστικά είναι 152 (τιμή) και 133 (φορτίο):
\mbox{153 στήλες $=$ 152 χαρακτηριστικά $+$ στόχος}, αντίστοιχα
\mbox{134 $=$ 133 $+$ στόχος} (πρβλ.\ Πίν.~\ref{tab:features_summary}).}
\end{table}

\begin{table}[htbp]
\centering
\caption{Ομάδες χαρακτηριστικών, αριθμός μεταβλητών ανά εργασία.}
\label{tab:features_summary}
\begin{tabular}{p{4cm}p{6.5cm}cc}
\toprule
Ομάδα & Περιγραφή & Τιμή & Φορτίο \\
\midrule
Ημερολογιακά & hour\_sin/cos, dow\_sin/cos, hour, dow, is\_holiday & 7 & 7 \\
Lags στόχου $y$ & $y_{t-k}$ για $k \in L$, y\_roll24, y\_roll168 & 10 & 10 \\
Lags φορτίου & $\text{load}_{t-k}$ για $k \in L$ & 8 & 8 \\
Πρόβλεψη φορτίου & load\_fc (proxy: load\_lag24) & 1 & --- \\
Lags υπολ.\ ζήτησης & residual\_load$_{t-k}$ για $k \in L$ & 8 & 8 \\
Lags ΑΠΕ παραγωγής & gen\_solar/wind$_{t-k}$ για $k \in L'$ & 10 & 10 \\
Τιμές καυσίμων & gas/co2 + lags $\times 2$ & 18 & --- \\
Μετεωρολογικά & 6 τοποθ.\ × 7 μεταβλ.\ + δείκτες απόντων & 84 & 84 \\
ENTSO-E ΑΠΕ forecast & solar/wind/gen\_fc + δείκτες απόντων & 6 & 6 \\
\midrule
\textbf{Σύνολο βασικό} & & \textbf{152} & \textbf{133} \\
Dense lags (MIMO/Direct) & $y_{t-k}$, $k \in \{4,5,7,\ldots,23\}$ & +18 & +18 \\
\textbf{Σύνολο dense} & & \textbf{170} & \textbf{151} \\
\bottomrule
\end{tabular}
\end{table}

Συνολικά, η διοχέτευση εναρμονίζει τις οκτώ ακατέργαστες πηγές σε ενιαία UTC-naive
ωριαία ανάλυση και τις διοχετεύει στις πέντε διαδοχικές φάσεις επεξεργασίας, παράγοντας
τους δύο πίνακες χαρακτηριστικών που αποθηκεύονται χωριστά σε μορφή \texttt{.parquet}
για τιμή και φορτίο. Το Σχήμα~\ref{fig:pipeline} αποτυπώνει αυτή τη ροή από τις πηγές
(αριστερά) στον τελικό πίνακα (δεξιά).

\begin{figure}[htbp]
\centering
\includegraphics[width=\textwidth]{fig_data_pipeline}
\caption{Διοχέτευση επεξεργασίας δεδομένων πέντε φάσεων: από τις οκτώ ακατέργαστες
πηγές (αριστερά) στους δύο τελικούς πίνακες χαρακτηριστικών (δεξιά).}
\label{fig:pipeline}
\end{figure}


\section[Κατασκευή Πίνακα Εκπαίδευσης]{Κατασκευή Πίνακα Εκπαίδευσης --- \texorpdfstring{\texttt{make\_xy()}}{make\_xy()}}
\label{sec:make_xy}

Πριν την εκπαίδευση κάθε μοντέλου, η \texttt{make\_xy()} αφαιρεί τον στόχο $y$ από τον
πίνακα, διατηρεί αποκλειστικά αριθμητικές στήλες, αντικαθιστά τιμές $\pm\infty$ με NaN
και εφαρμόζει \texttt{ffill()}. Τέλος αφαιρούνται γραμμές με NaN στον στόχο ή στα
χαρακτηριστικά.

Για τη Direct στρατηγική, η \texttt{\_make\_supervised\_xy()} κατασκευάζει πίνακα στόχου
$Y \in \mathbb{R}^{N \times H}$ με στήλες $[y_{t+1}, \ldots, y_{t+H}]$. Με το flag
\texttt{dense\_lags=True}, τα $y_{\text{lag1}}, \ldots, y_{\text{lag23}}$ προστίθενται στον
πίνακα $X$ πριν το fit.


\section{Χρονικός Διαχωρισμός Εκπαίδευσης / Αξιολόγησης}
\label{sec:train_test_split}

Ο διαχωρισμός εκπαίδευσης/αξιολόγησης ακολουθεί αυστηρά χρονολογική σειρά (expanding
window) χωρίς τυχαιοποίηση. Στην παρούσα εργασία εφαρμόζονται τέσσερα σχήματα
αξιολόγησης:

\begin{description}

  \item[1 Μήνας (Δεκέμβριος 2025).] Ίδιο σύνολο εκπαίδευσης, αξιολόγηση 744 ωρών.
  Επιτρέπει σύγκριση στρατηγικών ημερήσιας (24ωρης) έναντι εβδομαδιαίας (168ωρης)
  αρχικοποίησης.

  \item[Q1 2026 — Walk-forward μηνιαία επανεκπαίδευση.] Αξιολόγηση: Dec 1, 2025 -
  Feb 28, 2026 (2.160 ώρες). Ο κόφτης εκπαίδευσης κυλίεται μηνιαία:
  \{30/11/2025, 31/12/2025, 31/01/2026\}. Αξιολογεί αν η συχνή ενημέρωση οδηγεί σε
  σημαντική βελτίωση έναντι στατικά εκπαιδευμένων μοντέλων.
\end{description}

Η χρήση πολλαπλών σχημάτων αξιολόγησης επιτρέπει ανάλυση ευαισθησίας ως προς τον
ορίζοντα (βραχύς έναντι 3μηνου), τη σταθερότητα αγοράς (Δεκέμβριος έναντι Q4) και την
προσαρμοστικότητα (στατικά έναντι walk-forward).


\section{Περιγραφική Ανάλυση}
\label{sec:descriptive}

Η κατανόηση των στατιστικών ιδιοτήτων των δύο χρονοσειρών-στόχων είναι απαραίτητη
πριν την εκπαίδευση οποιουδήποτε μοντέλου. Η ανάλυση που ακολουθεί εξετάζει διαδοχικά
την πλήρη χρονοσειρά, την εποχική εξέλιξη ανά έτος, τα ωριαία και εβδομαδιαία
μοτίβα, τη στοχαστική δομή αυτοσυσχέτισης, τους χάρτες ενδοημερήσιας εξάρτησης,
και την αλληλεπίδραση τιμής με φυσικό αέριο και ΑΠΕ.


\subsection{Χρονοσειρά Τιμής Ηλεκτρικής Ενέργειας}
\label{subsec:price_timeseries}

Η χρονοσειρά καλύπτει περίπου εννέα χρόνια (Ιανουάριος 2017 - Νοέμβριος 2025), με
σύνολο $\sim$80.000 ωριαίων παρατηρήσεων. Το Σχήμα~\ref{fig:price_timeseries} αποτυπώνει
τη σειρά σε δύο επίπεδα ανάλυσης. Στο πάνω πάνελ φαίνεται η ωριαία τιμή (μπλε) μαζί με
τον 30-ήμερο κυλιόμενο μέσο (κόκκινη γραμμή), που αποκαλύπτει τις μεσοπρόθεσμες τάσεις
αφαιρώντας τη βραχυπρόθεσμη διακύμανση. Η σκιασμένη περιοχή Ιουνίου 2021 έως
Δεκεμβρίου 2022 σηματοδοτεί την ευρωπαϊκή ενεργειακή κρίση, κατά την οποία οι τιμές
υπερέβησαν τα 300~€/MWh. Από το 2023 και εξής παρατηρείται σταδιακή επιστροφή σε
προ-κρίσης επίπεδα. Το κάτω πάνελ παρουσιάζει τη μηνιαία κατανομή μέσω boxplots, που
αναδεικνύει την αυξημένη διασπορά των καλοκαιρινών μηνών (Ιούλιος-Αύγουστος) και
την εξαιρετική μεταβλητότητα των ετών κρίσης.

\begin{figure}[htbp]
  \centering
  \includegraphics[width=\textwidth]{fig_price_timeseries}
  \caption{Ωριαία τιμή ηλεκτρικής ενέργειας (€/MWh) στην ελληνική αγορά
           (Ιαν.\ 2017 - Μαρ.\ 2026). Πάνω: ωριαίες τιμές (μπλε) και 30-ήμερος
           κυλιόμενος μέσος (κόκκινος)· κίτρινη σκίαση για την περίοδο ενεργειακής
           κρίσης 2021-2022 και πράσινη για τα lockdown COVID-19. Κάτω: μηνιαία
           κατανομή τιμών (boxplots) για ολόκληρη την περίοδο.}
  \label{fig:price_timeseries}
\end{figure}

Στατιστικά, η τυπική απόκλιση της τιμής κυμαίνεται ανάλογα με το καθεστώς αγοράς:
33,6~€/MWh (Δεκέμβριος 2025, μέσος 110,4~€/MWh), 41,7~€/MWh (Ιανουάριος 2026) και
45,4~€/MWh (Φεβρουάριος 2026). Η πτώση του μέσου τον Φεβρουάριο 2026 (77,8~€/MWh έναντι
110~€/MWh τον Ιανουάριο) αποδίδεται σε αυξημένη υδροηλεκτρική παραγωγή και είναι το κεντρικό παράδειγμα μεταβολής καθεστώτος (regime shift) που αντιμετωπίζουν τα
μοντέλα στα πειράματα Q1~2026.

Σε αντίθεση, το φορτίο χαρακτηρίζεται από σαφώς μικρότερη σχετική μεταβλητότητα:
κυμαίνεται μεταξύ $\sim$4.500 και $\sim$9.000~MW με τυπική απόκλιση $\sim$18\% του μέσου,
έναντι $\sim$40\% για τιμή σε Q1~2026. Αυτή η διαφορά εξηγεί γιατί τα μοντέλα φορτίου
επιτυγχάνουν sMAPE $\sim$1,5-2\% ενώ τα μοντέλα τιμής κυμαίνονται στο $\sim$15-25\%.


\FloatBarrier
\subsection{Χρονοσειρά Ηλεκτρικού Φορτίου}
\label{subsec:load_timeseries}

Το Σχήμα~\ref{fig:load_timeseries} αποτυπώνει την αντίστοιχη χρονοσειρά φορτίου με την
ίδια διάταξη δύο πάνελ. Στο πάνω πάνελ η ωριαία κατανάλωση (πορτοκαλί) και ο 30-ήμερος
κυλιόμενος μέσος (κόκκινη γραμμή) αναδεικνύουν το επαναλαμβανόμενο εποχικό μοτίβο,
χωρίς τις απότομες μεταβολές καθεστώτος που χαρακτηρίζουν την τιμή. Το κάτω πάνελ
παρουσιάζει τη μηνιαία κατανομή μέσω boxplots, με σταθερότερη διασπορά διαχρονικά. Η
οπτική σύγκριση των δύο σειρών αναδεικνύει άμεσα τη μεγαλύτερη κανονικότητα του φορτίου,
που εξηγεί την ανώτερη προβλεψιμότητά του.

\begin{figure}[htbp]
  \centering
  \includegraphics[width=\textwidth]{fig_load_timeseries}
  \caption{Ωριαίο ηλεκτρικό φορτίο (MW) στην ελληνική αγορά (Ιαν.\ 2017 - Νοεμ.\ 2025).
           Πάνω: ωριαίο φορτίο (πορτοκαλί) και 30-ήμερος κυλιόμενος μέσος (κόκκινος)·
           πράσινη σκίαση για τα lockdown COVID-19. Κάτω: μηνιαία κατανομή φορτίου
           (boxplots) για ολόκληρη την περίοδο.}
  \label{fig:load_timeseries}
\end{figure}

Παρά τη φαινομενική ομαλότητα, η συγκεντρωτική εννεαετής απεικόνιση συμπιέζει τη
λεπτομερή εποχική δομή και των δύο σειρών: η μακροχρόνια μετατόπιση επιπέδου επικαλύπτει
το επαναλαμβανόμενο ετήσιο μοτίβο. Για να διαχωριστούν οι δύο αυτές συνιστώσες, η επόμενη
ενότητα σχεδιάζει τα μηνιαία προφίλ κάθε έτους χωριστά, αναδεικνύοντας ποιο μέρος της
διακύμανσης είναι σταθερά εποχικό και ποιο οφείλεται σε μεταβολές καθεστώτος.


\FloatBarrier
\subsection{Εποχιακή Εξέλιξη ανά Έτος}
\label{subsec:yearly_evolution}

Για να διαχωριστεί η βραχυπρόθεσμη από τη μακροπρόθεσμη εποχικότητα, το
Σχήμα~\ref{fig:monthly_avg_by_year} παρουσιάζει τη μέση μηνιαία τιμή ΗΑ και το μέσο
μηνιαίο φορτίο για κάθε έτος ξεχωριστά. Στις τιμές (πάνω πάνελ), τα έτη 2021 και 2022
(έντονες γραμμές) αποκλίνουν από την υπόλοιπη ομάδα: κορυφές >200~€/MWh
το καλοκαίρι του 2022 αντικατοπτρίζουν τα υψηλά επίπεδα TTF και τη μειωμένη
υδροηλεκτρική παραγωγή. Η ύπαρξη αυτών των «μεικτών καθεστώτων» (mixed regimes) στο
σύνολο εκπαίδευσης δυσκολεύει τη γενίκευση των μοντέλων σε νέες συνθήκες αγοράς.

Στο φορτίο (κάτω πάνελ), η εικόνα είναι σαφώς πιο σταθερή: το εποχικό μοτίβο (αιχμή
θέρους Ιούλιος-Αύγουστος, δευτερεύουσα αιχμή χειμώνα Ιανουάριος) επαναλαμβάνεται
ομοιόμορφα σε όλα τα έτη με μικρή χρονο-διακύμανση. Αυτή η σταθερότητα εξηγεί την
ανώτερη προβλεψιμότητα του φορτίου σε σχέση με την τιμή.

\begin{figure}[htbp]
  \centering
  \includegraphics[width=\textwidth]{fig_monthly_avg_by_year}
  \caption{Μηνιαία μέση τιμή ηλεκτρικής ενέργειας (πάνω) και φορτίου (κάτω)
           ανά έτος (2017-2025). Κάθε γραμμή αντιστοιχεί σε ένα ημερολογιακό έτος·
           τα έτη 2021-2022 σχεδιάζονται με έντονη γραμμή.}
  \label{fig:monthly_avg_by_year}
\end{figure}


\FloatBarrier
\subsection{Ωριαία και Εβδομαδιαία Εποχικότητα}
\label{subsec:seasonality}

Πέρα από τη μηνιαία εποχικότητα, η ημερήσια και εβδομαδιαία δομή αποτελούν τα κυρίαρχα
μοτίβα για τη βραχυπρόθεσμη πρόβλεψη. Το Σχήμα~\ref{fig:daily_weekly_profile}
αποτυπώνει το ενδοημερήσιο προφίλ τιμής σε τρία πάνελ: ανά εποχή, ανά ημέρα
εβδομάδας και ως θερμικό χάρτη (ημέρα $\times$ ώρα). Φαίνονται δύο χαρακτηριστικές
ωριαίες αιχμές: πρωινή (08:00-10:00, εκκίνηση βιομηχανίας) και βραδινή
(19:00-21:00, επιστροφή νοικοκυριών), με νυκτερινό βαθούλωμα (03:00-05:00). Το
εποχικό πάνελ αναδεικνύει την υψηλότερη στάθμη και μεταβλητότητα του χειμώνα και
του καλοκαιριού, ενώ το εβδομαδιαίο μοτίβο επιβεβαιώνει τη χαρακτηριστική πτώση
τιμής το Σαββατοκύριακο ($\sim$15\% κάτω από τον μέσο εργάσιμων ημερών), λόγω
μειωμένης βιομηχανικής ζήτησης.

\begin{figure}[htbp]
  \centering
  \includegraphics[width=\textwidth]{fig_daily_weekly_profile}
  \caption{Ενδοημερήσια μοτίβα τιμής ΗΑ. Πάνω: μέσο ωριαίο προφίλ ανά εποχή·
           κάτω αριστερά: μέσο ωριαίο προφίλ ανά ημέρα εβδομάδας (εργάσιμες μπλε,
           Σαββατοκύριακο πορτοκαλί)· κάτω δεξιά: θερμικός χάρτης μέσης τιμής (ημέρα
           εβδομάδας $\times$ ώρα).}
  \label{fig:daily_weekly_profile}
\end{figure}

Το Σχήμα~\ref{fig:daily_weekly_profile_load} παρουσιάζει το αντίστοιχο ανάλυση για
το φορτίο. Το ωριαίο προφίλ εμφανίζει σαφή αιχμή μεσημεριανής ζήτησης (12:00-13:00)
και στενότερη ζώνη $\pm 1\sigma$ σε σχέση με την τιμή, αντικατοπτρίζοντας τη μεγαλύτερη
ρυθμικότητα του φορτίου. Η εβδομαδιαία πτώση Σαββατοκύριακου υπάρχει αλλά είναι
λιγότερο έντονη. Αυτά τα μοτίβα δικαιολογούν άμεσα τη συμπερίληψη των ημερολογιακών
χαρακτηριστικών \texttt{hour\_sin/cos} και \texttt{dow\_sin/cos} στο σύνολο εισόδου
των μοντέλων.

\begin{figure}[htbp]
  \centering
  \includegraphics[width=\textwidth]{fig_daily_weekly_profile_load}
  \caption{Ενδοημερήσια μοτίβα ηλεκτρικού φορτίου, με ίδια διάταξη όπως η τιμή
           (Σχ.~\ref{fig:daily_weekly_profile}). Πάνω: προφίλ ανά εποχή·
           κάτω αριστερά: προφίλ ανά ημέρα εβδομάδας· κάτω δεξιά: θερμικός χάρτης μέσου
           φορτίου (ημέρα εβδομάδας $\times$ ώρα).}
  \label{fig:daily_weekly_profile_load}
\end{figure}


\FloatBarrier
\subsection{Αυτοσυσχέτιση και Επιλογή Υστερήσεων}
\label{subsec:autocorrelation}

Η συνάρτηση αυτοσυσχέτισης (ACF) και η μερική αυτοσυσχέτιση (PACF) ποσοτικοποιούν
τη στοχαστική δομή και τεκμηριώνουν την επιλογή του συνόλου υστερήσεων
$L$~\cite{weron2014electricity}. Χρησιμοποιούμε την ACF/PACF εδώ αποκλειστικά ως
εργαλείο επιλογής χαρακτηριστικών, όχι για βαθμονόμηση SARIMA.

Στο Σχήμα~\ref{fig:acf_pacf_price} η ACF της τιμής εμφανίζει κυματοειδή μορφή με
ισχυρές κορυφές στα lag-24, lag-48, $\ldots$, lag-168 (ημερήσιες αρμονικές στα
πολλαπλάσια των 24 ωρών και εβδομαδιαία αρμονική στο lag-168), που τεκμηριώνουν την
πολλαπλή εποχικότητα. Ξεχωρίζει η υψηλή στάθμη βάσης: ακόμη και στα ενδιάμεσα κοιλώματα η ACF
παραμένει κοντά στο $0{,}7$ και δεν αποσβένει έως το lag-168, ένδειξη μακράς μνήμης
που οφείλεται στα αργά μεταβαλλόμενα καθεστώτα τιμής (επίπεδα καυσίμων, ενεργειακή
κρίση). Η PACF κυριαρχείται από το lag-1 ($\approx 1{,}0$) με δεύτερη συνεισφορά στο
lag-2 ($\approx 0{,}8$) και ταχεία απόσβεση, με αμυδρές κορυφές γύρω από τα lag-24
και lag-48.

\begin{figure}[htbp]
  \centering
  \includegraphics[width=1.18\textwidth,center]{fig_acf_pacf_price_neo}
  \caption{Συνάρτηση αυτοσυσχέτισης (ACF, αριστερά) και μερικής αυτοσυσχέτισης
           (PACF, δεξιά) της ωριαίας τιμής ΗΑ για lags 1-168 ώρες. Οι κατακόρυφες
           διακεκομμένες γραμμές σημειώνουν τα lag-24, 48 και 168.}
  \label{fig:acf_pacf_price}
\end{figure}

Η μακρά μνήμη της τιμής αντιπαρατίθεται με την πολύ πιο τοπική εξάρτηση του φορτίου.
Το αντίστοιχο διάγραμμα για το φορτίο (Σχήμα~\ref{fig:acf_pacf_load}) παρουσιάζει
καθαρότερο ημερήσιο κύκλο: η ACF ταλαντώνεται έντονα μεταξύ βαθιών κοιλωμάτων
($\approx 0{,}1$ στα 12ωρα αντι-φάσης) και κορυφών $\approx 0{,}9$ στα πολλαπλάσια
των 24 ωρών. Η PACF κυριαρχείται και εδώ από το lag-1, με ασθενέστερη συνεισφορά
lag-2 ($\approx 0{,}5$) από ό,τι στην τιμή.

\begin{figure}[htbp]
  \centering
  \includegraphics[width=1.18\textwidth,center]{fig_acf_pacf_load_neo}
  \caption{Συνάρτηση αυτοσυσχέτισης (ACF, αριστερά) και μερικής αυτοσυσχέτισης
           (PACF, δεξιά) του ωριαίου ηλεκτρικού φορτίου για lags 1-168 ώρες. Οι
           κατακόρυφες διακεκομμένες γραμμές σημειώνουν τα lag-24, 48 και 168.}
  \label{fig:acf_pacf_load}
\end{figure}

Η αντιπαραβολή των δύο διαγραμμάτων αναδεικνύει τη βασική διαφορά των δύο σειρών.
Στην τιμή η ημερήσια ταλάντωση υπερτίθεται σε υψηλό υπόβαθρο επιμονής (κοιλώματα
$\approx 0{,}7$), ενώ στο φορτίο ο ημερήσιος κύκλος κυριαρχεί με κοιλώματα
που πέφτουν κοντά στο μηδέν· το φορτίο είναι πιο τακτικό και η δομή του πιο
προβλέψιμη. Και στις δύο σειρές η PACF επικεντρώνεται στο lag-1, δομή τύπου AR(1),
που μαζί με τις εποχικές κορυφές δικαιολογεί άμεσα το βασικό σύνολο
$L = \{1, 2, 3, 6, 12, 24, 48, 168\}$: καλύπτει τις στατιστικά σημαντικές υστερήσεις
χωρίς πλεονασμό.


\FloatBarrier
\subsection{Ενδοημερήσια Εξάρτηση Ημέρας D από D$-$1}
\label{subsec:heatmap}

Για την ποσοτικοποίηση της ισχυρής ενδοημερήσιας αυτοσυσχέτισης κατασκευάστηκαν
πίνακες Pearson $24 \times 24$, στους οποίους κάθε στοιχείο $(h, l)$ εκφράζει:

\begin{equation}
  \rho_{hl} = \mathrm{corr}\!\left(r_t^{(h)},\; r_{t-1}^{(l)}\right),
  \label{eq:intraday_corr}
\end{equation}

όπου $r_t^{(h)} = y_{t,h} - \bar{y}_t$ είναι η αποκεντρωμένη τιμή της ώρας $h$ ημέρας
$t$. Η αφαίρεση του ημερήσιου μέσου αναδεικνύει την καθαρή ενδοημερήσια δομή.

Στο Σχήμα~\ref{fig:corr_heatmap_price} για την τιμή, η κύρια διαγώνιος ($h = l$)
εμφανίζει τις υψηλότερες συσχετίσεις ($r > 0{,}8$): η τιμή της ίδιας ώρας την
προηγούμενη ημέρα είναι ο ισχυρότερος μεμονωμένος προβλεπτής. Οι ισχυρές
off-diagonal συσχετίσεις μεταξύ γειτονικών ωρών αιχμής (08:00-20:00) δικαιολογούν
την προσθήκη dense lags ($y_{t-4}, \ldots, y_{t-23}$) για τα MIMO/Direct μοντέλα.

\begin{figure}[htbp]
  \centering
  \includegraphics[width=\textwidth]{fig_corr_heatmap_price}
  \caption{Χάρτης Pearson ενδοημερήσιας συσχέτισης ($24 \times 24$) μεταξύ ώρας
           $h$ ημέρας $D$ και ώρας $l$ ημέρας $D-1$ για την τιμή ΗΑ (2017-2025).
           Οι ισοϋψείς γραμμές σημειώνουν τα $r=0{,}7$ και $r=0{,}9$.}
  \label{fig:corr_heatmap_price}
\end{figure}

Η ίδια αντίθεση τιμής-φορτίου που παρατηρήθηκε στην ACF επανεμφανίζεται και στο
επίπεδο της ενδοημερήσιας συσχέτισης μεταξύ διαδοχικών ημερών. Το φορτίο
(Σχήμα~\ref{fig:corr_heatmap_load}) εμφανίζει ισχυρότερη και πιο «καθαρή»
δομή ($\rho_{hh}$ έως $\sim 0{,}9$), συνεπές με τη μεγαλύτερη ημερήσια τακτικότητά
του. Η δομή κατά μπλοκ είναι πιο εμφανής σε σχέση με την τιμή, επιβεβαιώνοντας ότι
το lag-24 αποτελεί ακόμα ισχυρότερο χαρακτηριστικό για φορτίο από ό,τι για τιμή.

\begin{figure}[htbp]
  \centering
  \includegraphics[width=\textwidth]{fig_corr_heatmap_load}
  \caption{Χάρτης Pearson ενδοημερήσιας συσχέτισης ($24 \times 24$) μεταξύ ώρας
           $h$ ημέρας $D$ και ώρας $l$ ημέρας $D-1$ για το ηλεκτρικό
           φορτίο (2017-2025). Η κλίμακα χρωμάτων είναι ανεξάρτητη, προσαρμοσμένη
           στο υψηλότερο εύρος συσχέτισης του φορτίου (πρβλ.\ Σχ.~\ref{fig:corr_heatmap_price}).}
  \label{fig:corr_heatmap_load}
\end{figure}


\FloatBarrier
\subsection{Επίδραση Φυσικού Αερίου και Ανανεώσιμων Πηγών}
\label{subsec:fuel_res}

Στην ελληνική αγορά, οι μονάδες φυσικού αερίου καθορίζουν συχνά οριακά την τιμή
εκκαθάρισης. Αυτό αποτυπώνεται στο Σχήμα~\ref{fig:price_vs_gas}: η συσχέτιση τιμής
ΗΑ - τιμής TTF ανέρχεται σε $r \approx 0{,}74$ για την περίοδο 2017-2025. Η κλίση
της OLS αντικατοπτρίζει πόσα €/MWh αυξάνεται η τιμή για κάθε +1~€/MWh αύξηση του
αερίου. Η σημαντική διασπορά, ιδίως σε έτη υψηλής ΑΠΕ παραγωγής (2020, 2023), είναι
αναμενόμενη: όταν οι ανανεώσιμες πηγές εκτοπίζουν τις αεριοστροβιλικές μονάδες, η
σχέση αποσυνδέεται. Αυτή η υπό συνθήκη μη γραμμικότητα αποτελεί έναν από τους
βασικούς λόγους που τα μοντέλα ML υπερτερούν γραμμικών παλινδρομήσεων.

\begin{figure}[htbp]
  \centering
  \includegraphics[width=\textwidth]{fig_price_vs_gas}
  \caption{Αριστερά: διάγραμμα διασποράς ημερήσιας μέσης τιμής ΗΑ έναντι τιμής φυσικού
           αερίου TTF (2017-2025), χρωματισμένο ανά έτος, με προσαρμογή OLS (κόκκινη
           γραμμή) και αναγραφή της κλίσης της. Δεξιά: χρονοσειρά τιμής ΗΑ και τιμής
           αερίου με διπλό κατακόρυφο άξονα, που αναδεικνύει τη συγχρονισμένη εξέλιξή τους.}
  \label{fig:price_vs_gas}
\end{figure}

Συγκεντρωτικά, η περιγραφική ανάλυση σκιαγραφεί δύο σειρές με διακριτό χαρακτήρα: μια
τιμή με έντονη μεταβλητότητα, μακρά μνήμη και υπό συνθήκη εξάρτηση από τα καύσιμα, και
ένα φορτίο με ισχυρή ημερήσια τακτικότητα και υψηλή προβλεψιμότητα. Και οι δύο όμως
μοιράζονται την πολλαπλή εποχικότητα (ημερήσια, εβδομαδιαία, ετήσια) και τις ασυνέχειες
καθεστώτος που τεκμηρίωσαν τα προηγούμενα σχήματα. Αυτά τα χαρακτηριστικά καθορίζουν τόσο
την επιλογή χαρακτηριστικών όσο και τον τρόπο αξιολόγησης: κάθε σχήμα εκπαίδευσης πρέπει
να σέβεται τη χρονική σειρά και να εκθέτει τα μοντέλα στις ίδιες μεταβολές καθεστώτος που
θα αντιμετωπίσουν στην πράξη. Η επόμενη ενότητα οπτικοποιεί ακριβώς αυτόν τον διαχωρισμό.


\FloatBarrier
\subsection{Διαχωρισμός Εκπαίδευσης και Αξιολόγησης}
\label{subsec:train_test_viz}

Η αυστηρά χρονολογική διαίρεση (expanding window, χωρίς τυχαιοποίηση) αποτελεί
μεθοδολογική απαίτηση για EPF: οποιαδήποτε τυχαία κατανομή θα επέτρεπε στο μοντέλο
να «βλέπει» μελλοντικές τιμές κατά την εκπαίδευση, οδηγώντας σε φαινομενικά άριστα
αλλά μη αναπαραγώγιμα αποτελέσματα~\cite{bergmeir2018}. Το
Σχήμα~\ref{fig:train_test_split} οπτικοποιεί τον διαχωρισμό για την περίοδο
αξιολόγησης Q1~2026 (Δεκ.\ 2025 - Φεβ.\ 2026). Η εκπαίδευση καλύπτει ολόκληρο το
ιστορικό (Ιαν.\ 2017 - Νοεμ.\ 2025). Η αξιολόγηση γίνεται σε walk-forward σχήμα: για
το Q1~2026 ο κόφτης εκπαίδευσης κυλίεται μηνιαία (\{30/11/2025, 31/12/2025,
31/01/2026\}), χωρίς validation set εντός της περιόδου αξιολόγησης. Το σχήμα
παρουσιάζει δύο πάνελ (τιμή πάνω, φορτίο κάτω)· η κατακόρυφη διακεκομμένη γραμμή
σημειώνει το σημείο κοπής εκπαίδευσης (τέλος Νοεμβρίου 2025) και η σκιασμένη περιοχή
την περίοδο δοκιμής.

\begin{figure}[htbp]
  \centering
  \includegraphics[width=\textwidth]{fig_train_test_split}
  \caption{Χρονικός διαχωρισμός εκπαίδευσης (Ιαν.\ 2017 - Νοεμ.\ 2025) και
           αξιολόγησης (κόκκινο, Δεκ.\ 2025 - Φεβ.\ 2026), για την τιμή (πάνω) και το
           φορτίο (κάτω).}
  \label{fig:train_test_split}
\end{figure}
